\subsection{Еще симметричные и бесконечные цепи}
%%%%%%%%%%%%%%%%%%%%%%%%%%%%%%%%%%%%%%%%%%%%%%%%%%%%
%%%%%%%%%%%%%%%%%%%%%%%%%%%%%%%%%%%%%%%%%%%%%%%%%%%%1
 \z{Определите эквивалентные сопротивления бесконечной цепи из резисторов, схема которой приведены на рисунке.}
%
{$R_\text{э}=\frac{3+\sqrt{21}}{2}R$}

%%%%%%%%%%%%%%%%%%%%%%%%%%%%%%%%%%%%%%%%%%%%%%%%%%%%1
 \z{Определите эквивалентные сопротивления бесконечной цепи из резисторов, схема которой приведены на рисунке.}
%
{$R_\text{э}=\frac{1+\sqrt{5}}{2}R$}

%%%%%%%%%%%%%%%%%%%%%%%%%%%%%%%%%%%%%%%%%%%%%%%%%%%%1
 \z{Очень длинная цепочка составлена из одинаковых батареек, имеющих ЭДС $\emf= 12\, \text{В}$ и внутренние сопротивления $r = 4\, \text{Ом}$, и резисторов с сопротивлением $R = 15\, \text{Ом}$ (см. рисунок). Определите ЭДС $\emf_0$ и внутреннее сопротивление $r_0$ эквивалентной батарейки.
\begin{figure}[!h]
    \centering
    \begin{minipage}[c]{0.32\linewidth} % [c] вместо [t]
        \centering
        \includegraphics[height=2cm]{3.9.1} % фиксированная высота
        \caption*{Задача 1}
    \end{minipage}
    \hfill
    \begin{minipage}[c]{0.32\linewidth}
        \centering
        \includegraphics[height=2cm]{3.9.2}
        \caption*{Задача 2}
    \end{minipage}
    \hfill
    \begin{minipage}[c]{0.32\linewidth}
        \centering
        \includegraphics[height=2cm]{3.9.3}
        \caption*{Задача 3}
    \end{minipage}
\end{figure}
 }
%
{$\emf_0=30\,\text{В};~~~r_0=10\,\text{Ом}$}

%%%%%%%%%%%%%%%%%%%%%%%%%%%%%%%%%%%%%%%%%%%%%%%%%%%%1
 \z{Найдите сопротивление между клеммами $A$ и $B$ бесконечной цепи, схема которой изображена на рисунке. Сопротивление каждого резистора равно $R$.}
%
{$R_{AB}=\frac{2\sqrt{3}}{3}R$}

%%%%%%%%%%%%%%%%%%%%%%%%%%%%%%%%%%%%%%%%%%%%%%%%%%%%1
 \z{Найдите сопротивление между клеммами $A$ и $B$ бесконечной цепи, схема которой изображена на рисунке. Сопротивление каждого резистора равно $R$.}
%
{$R_{AB}=\sqrt{2}R$}

%%%%%%%%%%%%%%%%%%%%%%%%%%%%%%%%%%%%%%%%%%%%%%%%%%%%1
 \z{Найдите сопротивление цепи между точками $A$ и $B$, если сопротивление отрезка провода между $A$ и $D$ равно $R$, все фигуры на рисунке являются квадратами, которых бесконечно много, и они соединены в серединах сторон.
\begin{figure}[!h]
    \centering
    \begin{minipage}[c]{0.32\linewidth} % [c] вместо [t]
        \centering
        \includegraphics[height=3cm]{3.9.4} % фиксированная высота
        \caption*{Задача 4}
    \end{minipage}
    \hfill
    \begin{minipage}[c]{0.32\linewidth}
        \centering
        \includegraphics[height=3cm]{3.9.5}
        \caption*{Задача 5}
    \end{minipage}
    \hfill
    \begin{minipage}[c]{0.32\linewidth}
        \centering
        \includegraphics[height=3cm]{3.9.6}
        \caption*{Задача 6}
    \end{minipage}
\end{figure}
 }
%
{$R_{AB}=1.32R$}

%%%%%%%%%%%%%%%%%%%%%%%%%%%%%%%%%%%%%%%%%%%%%%%%%%%%1
 \z[{\begin{wrapfigure}[3]{r}{0.3\textwidth}%
		\vspace{-15pt}%
		\begin{flushright}%
 			\includegraphics[width=0.3\textwidth]{3.9.7}%
		\end{flushright}%
	\end{wrapfigure}%
	}]%
 {Определите сопротивление бесконечной сетки с сопротивлением одной стороны ячейки $r$, измеренное между узлами $A$ и $B$.}
%
{$R_{AB}=\frac{r}{3}$}




